\documentclass[10pt]{article}

\usepackage[margin=0.5in,top=0.35in,bottom=0.4in]{geometry}
\usepackage[T1]{fontenc}
\usepackage{lmodern}
\usepackage{microtype}
\usepackage{amssymb}
\usepackage{xcolor}
\usepackage{titlesec}
\usepackage{enumitem}
\usepackage[hidelinks]{hyperref}

\definecolor{accent}{HTML}{1F4E79}
\definecolor{contactgray}{HTML}{444444}

\pagestyle{empty}
\setlength{\parindent}{0pt}
\linespread{0.95}

\titleformat{\section}
  {\large\bfseries\color{accent}}{}{0pt}{}
  [\vspace{1pt}\textcolor{accent}{\rule{\textwidth}{0.6pt}}]
\titlespacing*{\section}{0pt}{6pt}{2.5pt}

\setlist[itemize]{leftmargin=1.25em,itemsep=0.2pt,topsep=1pt,parsep=0pt}

\newcommand{\projtitle}[1]{{\fontsize{9.5}{11.4}\selectfont #1}}

% Change for each target vacancy.
\newcommand{\TargetTitle}{Senior Scientific ML Researcher \,|\, AI for Science \& ML Systems}
\newcommand{\TargetFocus}{Self-supervised learning $\cdot$ representation learning $\cdot$ PyTorch $\cdot$ GPU/HPC $\cdot$ multiscale physical fields}

\begin{document}

% ---------- Header ----------
\begin{center}
  {\Huge\bfseries Guang-Xing Li, Ph.D.}\\[1.5pt]
  {\large\color{accent}\bfseries \TargetTitle}\\[1.5pt]
  {\small\color{contactgray} \TargetFocus}\\[2pt]
  {\small\color{contactgray}
   \href{mailto:ligx.ngc7293@gmail.com}{ligx.ngc7293@gmail.com} \quad$|$\quad
   \href{https://github.com/gxli}{github.com/gxli} \quad$|$\quad
   \href{https://gxli.github.io}{gxli.github.io}}
\end{center}
\vspace{1.5pt}

% ---------- Technical Summary ----------
\section*{Technical Summary}
Computational physicist and ML engineer building scientific-ML systems for high-dimensional data,
including latent predictive models, from representation architecture and PyTorch implementation
through GPU training, dense inference, validation, and open-source release. Creator of ScaleAware-JEPA, a world-model-inspired JEPA system for
discovering structure in continuous physical fields.

% ---------- Technical Skills (top) ----------
\section*{Technical Skills}
\begin{itemize}
  \item \textbf{ML/AI:} Python, PyTorch, self-supervised learning, JEPA, ConvNeXt, representation
        learning, PCA/UMAP
  \item \textbf{ML Systems:} GPU training, gradient accumulation, checkpointing, tiled inference,
        config-driven experiments, testing
  \item \textbf{Scientific Computing:} NumPy, SciPy, HPC/Linux, large-scale simulation,
        high-dimensional scientific data
  \item \textbf{Software:} Git, open-source packaging, reproducible pipelines, documentation
\end{itemize}

% ---------- Selected Technical Projects ----------
\section*{Flagship ML System}

\textbf{ScaleAware-JEPA} --- \textit{Latent predictive model for scientific discovery}
\hfill{\footnotesize\color{contactgray!70} \textbf{End-to-end developer} $\cdot$ 2026 $\cdot$ Sole author}\\[2pt]
{\fontsize{6.5}{7.8}\selectfont\color{contactgray!50}\itshape Paper: ScaleAware-JEPA: Latent Representation for
Discovery in Multiscale Physical Fields $\cdot$
\href{https://arxiv.org/abs/2606.29723}{arXiv:2606.29723}\\
Code: \href{https://github.com/gxli/SA-JEPA}{github.com/gxli/SA-JEPA}}
\begin{itemize}
  \item \textbf{Problem:} discover meaningful structure in complex physical data \textbf{without
        predefined labels, objects, or segmentation rules}.
  \item \textbf{Architecture:} predicts hidden structure \textbf{in latent space rather than
        reconstructing raw data}; CDD provides pixel-registered scale components and
        physical-scale coordinates, while a ConvNeXt V2-style encoder fuses information across
        scales and produces \textbf{dense 32-D latent fields}. Scale-aware masking adapts the
        prediction context to the intrinsic scale hierarchy rather than fixed image patches.
  \item \textbf{Engineering:} config-driven PyTorch training with gradient accumulation,
        checkpointing, tiled inference, PCA/UMAP diagnostics, tests, scripts, documentation, and
        reproducible experiment configs.
  \item \textbf{Validation:} evaluated on a \textbf{$256\times 256\times 256$ MHD turbulence cube},
        molecular-gas observations, and urban nighttime-light data; learned dense structural atlases
        recover coherent morphology and \textbf{physically distinct scaling relations} without
        predefined labels or segmentation.
\end{itemize}

\section*{Technical Ownership \& Collaboration}
\begin{itemize}
  \item \textbf{Own SA-JEPA end to end:} architecture, PyTorch implementation, experiment design,
        validation, testing, documentation, release.
  \item Maintain open-source scientific-ML codebases with reproducible configs, tests, examples,
        and documentation.
\end{itemize}

\projtitle{\textbf{Scale-Aware Adversarial Analysis} --- \textit{Generative-model diagnostics}}
\hfill{\footnotesize\color{contactgray!70} 2026 $\cdot$ Co-author}\\[2pt]
{\fontsize{6.5}{7.8}\selectfont\color{contactgray!50}\itshape Paper: Scale-Aware Adversarial Analysis: A
Diagnostic for Generative AI in Multiscale Complex Systems $\cdot$
\href{https://arxiv.org/abs/2605.00510}{arXiv:2605.00510}}
\begin{itemize}
  \item \textbf{Method:} deterministic, physically constrained interventions in CDD scale space for
        evaluating DDPMs of multiscale systems.
  \item \textbf{Result:} exposed localized structural freezing and nonlinear instability missed by
        aggregate image-quality metrics.
\end{itemize}

\newpage

\section*{Additional Technical Projects}

\projtitle{\textbf{Adjacent Correlation Analysis (ACA)} --- \textit{Discovery of hidden regularities}}
\hfill{\footnotesize\color{contactgray!70} 2025 $\cdot$ Sole author}\\[2pt]
{\fontsize{6.5}{7.8}\selectfont\color{contactgray!50}\itshape Paper: Revealing Hidden Correlations from Complex
Spatial Distributions: Adjacent Correlation Analysis $\cdot$
\href{https://arxiv.org/abs/2506.05759}{arXiv:2506.05759}\\
Code: \href{https://github.com/gxli/Adjacent-Correlation-Analysis}{github.com/gxli/Adjacent-Correlation-Analysis}}
\begin{itemize}
  \item \textbf{Method:} recovers local relationships hidden by global statistics, representing
        complex spatial patterns as phase-space vector fields.
  \item \textbf{Result:} recovered vectors form the vector field on an attracting manifold---
        enabling classification, prediction, parameter fitting, and forecasting.
\end{itemize}

\projtitle{\textbf{Constrained Diffusion Decomposition (CDD)} --- \textit{Multiscale decomposition for
continuous physical data}}
\hfill{\footnotesize\color{contactgray!70} 2022 $\cdot$ Sole author}\\[2pt]
{\fontsize{6.5}{7.8}\selectfont\color{contactgray!50}\itshape Paper: Multi-scale Decomposition of Astronomical Maps:
A Constrained Diffusion Method $\cdot$ ApJS 259, 59 $\cdot$ \href{https://arxiv.org/abs/2201.05484}{arXiv:2201.05484}\\
Code: \href{https://github.com/gxli/Constrained-Diffusion-Decomposition}{github.com/gxli/Constrained-Diffusion-Decomposition} $\cdot$ PyPI: constrained-diffusion}
\begin{itemize}
  \item \textbf{Method:} nonlinear-diffusion decomposition producing pixel-registered physical
        scales without the ringing and negative artifacts introduced by conventional band-limited
        filtering.
  \item \textbf{Reuse:} multiscale representation used as an input coordinate for SA-JEPA and
        scale-aware generative-model diagnostics.
\end{itemize}

% ---------- Experience ----------
\section*{Experience}
\textbf{Principal Investigator \& Associate Professor} \hfill \textbf{2018--Present}\\
\textit{Yunnan University (SWIFAR), China}\\
\begin{itemize}
  \item Lead computational research projects spanning scientific ML, multiscale modeling, and
        high-dimensional observational and simulation data.
  \item Supervise PhD/MSc researchers in computational methods, scientific data analysis, and
        research software development.
  \item Advise collaborators on complex observational and simulation datasets, translating
        scientific questions into computational formulations, physical diagnostics, data pipelines,
        and ML workflows.
\end{itemize}
\vspace{3pt}

\textbf{DFG Research Fellow} \hfill \textbf{2014--2018}\\
\textit{LMU Munich (Universit\"ats-Sternwarte), Germany}\\
\begin{itemize}
  \item Led a DFG-funded program developing computational methods to extract physical understanding
        from observational and simulation data in Linux/HPC environments.
  \item Developed G-virial, combining analytical physics with numerical simulations for structure
        identification in turbulent data.
\end{itemize}
\vspace{3pt}

\textbf{Ph.D. Researcher} \hfill \textbf{2010--2014}\\
\textit{Max Planck Institute for Radio Astronomy, Germany}\\
\begin{itemize}
  \item Built scalable algorithms for multiscale structure extraction from high-dimensional,
        heterogeneous datasets.
\end{itemize}
\vspace{5pt}

% ---------- Academic Credibility ----------
\section*{Research Impact}
{\small \textbf{111 scientific publications} $\cdot$ \textbf{1,800+ citations} $\cdot$
\textbf{h-index 24} $\cdot$ \textbf{\texteuro280k+ PI-led competitive funding}}

% ---------- Education ----------
\vspace{4pt}
\section*{Education}
\textbf{Ph.D. Astrophysics / Computational Modeling} --- Max Planck Institute for Radio Astronomy \&
University of Bonn (2014), \textit{magna cum laude}\\
\textbf{M.Sc. Astrophysics} --- University of Science and Technology of China (2010)\\
\textbf{B.Sc. Electronics} --- Peking University (2007)

\end{document}
